\section{Conclusion}
\label{sec:conclusion}

This paper has presented a systematic literature review of guardrails for large
language models, focusing on two deployment-critical but underexplored dimensions:
real-time operation and multimodal contexts.

The central empirical finding is a pervasive reproducibility crisis.
Across 75 primary sources and 237 extracted claims, exactly five
claims survived author-verified scrutiny (three temporally separated self-review
passes by the same single author, each at least two weeks apart, with unanimous
agreement required; see Section~\ref{subsec:verification_protocol}).
The five confirmed findings are narrow but important: (1) the dual reward model
architecture and threshold arbitration from Llama~2~\cite{touvron2023llama2};
(2) safety context distillation as a weight-embedding technique~\cite{touvron2023llama2};
(3) the instruction-following classification framing of Llama Guard~\cite{inan2023llama};
(4) the multi-turn conversational priming attack amplifying jailbreaks by
29--71\%~\cite{russinovich2024greatescalation}; and (5) a meta-finding that
evaluation-protocol choices alone shift reported attack-success rates by up to
30 percentage points~\cite{li2024salad} --- an evaluation-fragility result that
is itself a cross-cutting finding from the 72 remaining corpus papers rather than
a result of any single study.
Everything else---the specific latency figures, the precise attack success rates,
the comparative accuracy claims---collapsed under structured self-review.

The taxonomy developed in this review organises existing mechanisms along three
axes---placement, implementation layer, and modality coverage---and reveals that:
real-time in-generation (streaming) safety checking has no verified solution;
audio and video safety is entirely absent from the academic literature despite
active deployment; and session-level context, necessary to resist multi-turn
attacks, is not incorporated in any verified runtime guardrail system.

The open-problems catalogue in Section~\ref{sec:open_problems} identifies nine
specific research gaps (RP1, RT1--RT3, MM1--MM3, FT1, FG1) across reproducibility,
streaming safety, cross-modal alignment, fine-tuning robustness, and formal
guarantees.
These represent the highest-priority directions for a field whose deployment
urgency has outpaced its scientific rigour.

We hope this review serves both as a reference for researchers seeking to build
on verified foundations and as a call to action for the evaluation-first methodology
that the field requires.
The per-claim audit artefact, the full annotated reference corpus (including the
complete per-paper modality matrix), and the reproducibility analysis code are
available from the corresponding author on request.

\section*{Acknowledgements}
The author thanks the open-source research community for maintaining publicly
accessible preprint repositories without which this review would not have been
possible.
Two discrepant inclusion decisions arising during the title/abstract screening
phase were resolved by consultation with external domain experts (one specialist
in LLM safety evaluation and one specialist in multimodal machine learning);
the author is grateful for their judgement on those borderline cases, which
affected the final corpus composition.
